\section{Fase 1 -- Preparación de infraestructura}

\subsection{Objetivo}
El laboratorio consiste en tres máquinas virtuales Ubuntu 22.04 orquestadas con Vagrant + libvirt sobre una red privada propia (\texttt{red-prueba}). El despliegue es totalmente reproducible mediante \emph{infraestructura como código}: un único \texttt{vagrant up} desde el host crea la red, los discos persistentes, las VMs, los usuarios personalizados y deja las máquinas listas para que cualquiera de los \emph{provisioners} de fases posteriores (Kubespray, monitoring, seguridad) pueda ejecutarse a demanda.

\subsection{Requisitos de la infraestructura}

A continuación se muestran las especificaciones técnicas asignadas a cada nodo en el \texttt{Vagrantfile}:

\begin{center}
  \begin{tabular}{|c|c|}
    \hline
    \multicolumn{2}{|c|}{Infraestructura nodo \texttt{test-master}} \\
    \hline
    Componente & Especificación \\
    \hline
    vCPU & 2 núcleos \\
    RAM & 4096 MB \\
    Disco principal & 40 GB (caja Vagrant) \\
    Disco persistente & 40 GB qcow2 (\texttt{/dev/vdb}) \\
    Red privada & \texttt{red-prueba} (NAT, \texttt{virbr2}) \\
    IP & 10.0.1.10 \\
    Sistema Operativo & Ubuntu 22.04 LTS \\
    Usuario personalizado & \texttt{master} \\
    \hline
  \end{tabular}
\end{center}

\begin{center}
  \begin{tabular}{|c|c|}
    \hline
    \multicolumn{2}{|c|}{Infraestructura nodo \texttt{test-worker1}} \\
    \hline
    Componente & Especificación \\
    \hline
    vCPU & 2 núcleos \\
    RAM & 3072 MB \\
    Disco principal & 40 GB (caja Vagrant) \\
    Disco persistente & 20 GB qcow2 (\texttt{/dev/vdb}) \\
    Red privada & \texttt{red-prueba} (NAT, \texttt{virbr2}) \\
    IP & 10.0.1.11 \\
    Sistema Operativo & Ubuntu 22.04 LTS \\
    Usuario personalizado & \texttt{worker1} \\
    \hline
  \end{tabular}
\end{center}

\begin{center}
  \begin{tabular}{|c|c|}
    \hline
    \multicolumn{2}{|c|}{Infraestructura nodo \texttt{test-worker2}} \\
    \hline
    Componente & Especificación \\
    \hline
    vCPU & 1 núcleo \\
    RAM & 2048 MB \\
    Disco principal & 40 GB (caja Vagrant) \\
    Disco persistente & 20 GB qcow2 (\texttt{/dev/vdb}) \\
    Red privada & \texttt{red-prueba} (NAT, \texttt{virbr2}) \\
    IP & 10.0.1.12 \\
    Sistema Operativo & Ubuntu 22.04 LTS \\
    Usuario personalizado & \texttt{worker2} \\
    \hline
  \end{tabular}
\end{center}

El reparto de RAM no es uniforme deliberadamente: \texttt{test-worker1} aloja los tres CMS vulnerables (DVWA, WebGoat, Juice Shop) junto a Suricata, mientras que \texttt{test-worker2} concentra todo el \emph{stack} de monitorización (Grafana, Prometheus, Loki) sobre \texttt{hostPath} local. \texttt{test-master} solo aloja el \emph{control-plane} de Kubernetes y etcd.

\subsection{Instalación de paquetería en el \emph{host}}

A continuación se listan los paquetes\footnote{El \emph{host} desde el que se desarrolla el proyecto es Arch Linux con \texttt{pacman} como gestor de paquetes. La reproducción es posible en otras distribuciones cambiando el gestor; los nombres de los paquetes son comunes en la mayoría de distros mainstream.} necesarios en la máquina anfitriona:

\begin{itemize}
  \item \textbf{Libvirt y QEMU}
    \begin{lstlisting}[style=bash]
sudo pacman -S qemu-full virt-manager virt-viewer spice-gtk virt-install libvirt
sudo pacman -S --needed dnsmasq vde2 iproute2 iptables-nft openbsd-netcat
    \end{lstlisting}

  \item \textbf{Activación de los servicios}
    \begin{lstlisting}[style=bash]
sudo systemctl enable --now libvirtd.service
sudo systemctl enable --now virtlogd.service
    \end{lstlisting}

  \item \textbf{Grupo libvirt}: necesario para ejecutar \texttt{virsh} sin privilegios de superusuario:
    \begin{lstlisting}[style=bash]
sudo gpasswd -a $USER libvirt
newgrp libvirt
    \end{lstlisting}

  \item \textbf{Vagrant} (vía AUR con \texttt{yay}):
    \begin{lstlisting}[style=bash]
yay -S vagrant
vagrant plugin install vagrant-libvirt
    \end{lstlisting}

  \item \textbf{Herramientas auxiliares} usadas por los \emph{scripts} del proyecto:
    \begin{itemize}
      \item \texttt{qemu-img}: creación de discos qcow2 (incluido en \texttt{qemu-full}).
      \item \texttt{setfacl}: gestión de ACL para permisos del usuario \texttt{libvirt-qemu} (paquete \texttt{acl}).
      \item \texttt{virsh}: gestión de dominios y redes de libvirt.
    \end{itemize}
\end{itemize}

\subsection{Definición de la red privada}

La red privada \texttt{red-prueba} es gestionada por libvirt y su definición vive en \texttt{network/red-prueba.xml}, dentro del propio repositorio. La automatización es responsabilidad de \texttt{scripts/network-check.sh}, que se ejecuta como \emph{trigger} \texttt{before :up} antes de levantar la primera VM.

\begin{lstlisting}[style=xml]
<network>
  <name>red-prueba</name>
  <forward mode='nat'/>
  <bridge name='virbr2' stp='on' delay='0'/>
  <ip address='10.0.1.1' netmask='255.255.255.0'>
    <dhcp>
      <range start='10.0.1.2' end='10.0.1.254'/>
      <host mac='52:54:00:00:01:10' name='test_master'  ip='10.0.1.10'/>
      <host mac='52:54:00:00:01:11' name='test_worker1' ip='10.0.1.11'/>
      <host mac='52:54:00:00:01:12' name='test_worker2' ip='10.0.1.12'/>
    </dhcp>
  </ip>
</network>
\end{lstlisting}

La asignación de IP estática por reserva DHCP \emph{con MAC fija} es obligatoria: si no se fijan las MAC, Vagrant crearía una segunda red automática (\emph{management}) y se perdería la trazabilidad de las IPs reales. La asignación final queda:

\begin{center}
  \begin{tabular}{|c|c|c|}
    \hline
    Nodo & IP & MAC \\
    \hline
    test-master  & 10.0.1.10 & 52:54:00:00:01:10 \\
    test-worker1 & 10.0.1.11 & 52:54:00:00:01:11 \\
    test-worker2 & 10.0.1.12 & 52:54:00:00:01:12 \\
    \hline
  \end{tabular}
\end{center}

\paragraph{Comprobación previa al \texttt{up}.} El \emph{script} \texttt{network-check.sh} no solo verifica la existencia de la red: compara la configuración \emph{activa} en libvirt con el XML del repositorio (ignorando UUIDs y MACs autogenerados) y, si difieren, destruye y redefine la red. Esto evita un fallo silencioso muy frecuente: tener una red \texttt{red-prueba} \emph{antigua} con un \emph{range} o un MAC distinto que no matchea la del \texttt{Vagrantfile} y descubrir solo más tarde que las VMs reciben IPs aleatorias.

\subsection{Vagrantfile}

El despliegue se realiza con Vagrant + el provider \texttt{vagrant-libvirt}. Se descartó Terraform por estar enfocado a infraestructura \emph{cloud} y no ofrecer integración directa con libvirt local sin un \emph{provider} de la comunidad poco mantenido. A continuación, el \texttt{Vagrantfile} completo:

\begin{lstlisting}[style=bash]
ENV['VAGRANT_NO_PARALLEL'] = 'yes'

Vagrant.configure("2") do |config|
  config.vm.box = "generic/ubuntu2204"

  # Synced folder: rsync (el unico modo fiable con vagrant-libvirt sin NFS).
  # Excluimos .git/.vagrant/disks para no rsynchronizar GB de qcow2.
  config.vm.synced_folder ".", "/vagrant", type: "rsync",
    rsync__exclude: [".git/", ".vagrant/", "disks/"]

  nodes = {
    "test-master"  => { "ip" => "10.0.1.10", "mac" => "52:54:00:00:01:10", "user" => "master"  },
    "test-worker1" => { "ip" => "10.0.1.11", "mac" => "52:54:00:00:01:11", "user" => "worker1" },
    "test-worker2" => { "ip" => "10.0.1.12", "mac" => "52:54:00:00:01:12", "user" => "worker2" }
  }

  nodes.each do |name, info|
    config.vm.define name do |node|
      node.vm.hostname = name

      node.vm.network "private_network",
                      ip: info["ip"],
                      netmask: "255.255.255.0",
                      mac: info["mac"],
                      libvirt__network_name: "red-prueba",
                      libvirt__forward_mode: "nat",
                      libvirt__always_destroy: false

      node.vm.provider :libvirt do |libvirt|
        libvirt.memory = case name
                         when "test-master" then 4096
                         when "test-worker1" then 3072
                         else 2048
                         end
        libvirt.cpus = (name == "test-master" || name == "test-worker1") ? 2 : 1
        libvirt.mgmt_attach = false  # sin red de management adicional
      end

      # Provisioners encadenados en todas las VMs:
      node.vm.provision "shell" do |s|
        s.path = "scripts/config.sh"        # paquetes base (vim, python3, git...)
      end
      node.vm.provision "shell" do |s|
        s.path = "scripts/ssh-setup.sh"     # crea el usuario custom + claves inter-nodo
        s.args = [info["user"]]
      end

      # Provisioners "run: never" solo en master, disparados manualmente:
      if name == "test-master"
        node.vm.provision "setup-kubespray", type: "shell", run: "never" do |s|
          s.path = "scripts/setup-kubespray.sh"
          s.privileged = false
        end
        node.vm.provision "setup-lab", type: "shell", run: "never" do |s|
          s.path = "scripts/setup-lab.sh"
          s.privileged = false
        end
        node.vm.provision "setup-security", type: "shell", run: "never" do |s|
          s.path = "scripts/setup-security.sh"
          s.privileged = false
        end
      end
    end
  end

  # Triggers del host (fuera del ciclo de provisioners):
  config.trigger.before :up, only_on: "test-master" do |trigger|
    trigger.name = "[+] Comprobacion de la red puente"
    trigger.run = { path: "scripts/network-check.sh" }
  end
  config.trigger.before :up, only_on: "test-master" do |trigger|
    trigger.name = "[+] Comprobacion del almacenamiento: discos qcow2"
    trigger.run = { path: "scripts/create-disks.sh" }
  end
  config.trigger.after :up, only_on: "test-worker2" do |trigger|
    trigger.name = "[+] Configuracion SSH las VMs"
    trigger.run = { path: "scripts/ssh-config.sh" }
  end
  config.trigger.after :up, only_on: "test-worker2" do |trigger|
    trigger.name = "[+] Conexion y montaje persistente de los discos"
    trigger.run = { path: "scripts/attach-disks.sh" }
  end
end
\end{lstlisting}

\paragraph{Despliegue secuencial.} La variable de entorno \texttt{VAGRANT\_NO\_PARALLEL = 'yes'} fuerza el arranque secuencial de las VMs. Esto evita condiciones de carrera entre \emph{triggers} y \emph{provisioners}: por ejemplo, \texttt{ssh-config.sh} se ejecuta como \emph{trigger after :up, only\_on: test-worker2} y necesita que las tres VMs estén ya arrancadas para escribir las tres entradas en \texttt{\textasciitilde/.ssh/config} del \emph{host}. Sin \texttt{NO\_PARALLEL}, Vagrant podría lanzar el \emph{trigger} antes de que \texttt{vagrant ssh-config} devolviera info de los tres nodos.

\paragraph{Provisioners \texttt{run: "never"}.} Las fases pesadas (\texttt{setup-kubespray}, \texttt{setup-lab}, \texttt{setup-security}) están marcadas como \texttt{run: "never"}: no se ejecutan en \texttt{vagrant up}. Se invocan manualmente:

\begin{lstlisting}[style=bash]
vagrant provision test-master --provision-with setup-kubespray
vagrant provision test-master --provision-with setup-lab
vagrant provision test-master --provision-with setup-security
\end{lstlisting}

Esto permite reconstruir la infraestructura (\texttt{vagrant destroy -f \&\& vagrant up}) en pocos minutos sin pagar el coste de re-instalar Kubespray (~10\,min) o re-desplegar el \emph{stack} de monitoring. Cada \emph{provisioner} es idempotente, así que puede re-ejecutarse sin romper estado.

\subsection{Provisioners de fase 0 (\texttt{config.sh} y \texttt{ssh-setup.sh})}

Dos \emph{provisioners shell} corren automáticamente en cada \texttt{up} dentro de las tres VMs.

\subsubsection{\texttt{config.sh}}

Instala los paquetes base que necesitan las fases posteriores (\texttt{python3} y \texttt{python3-venv} para Kubespray, \texttt{git} para clonar repositorios, \texttt{vim} para edición manual). Es deliberadamente minimalista: Kubespray y Helm tienen sus propios bootstrap.

\subsubsection{\texttt{ssh-setup.sh}}

Crea un usuario personalizado por VM (\texttt{master}, \texttt{worker1}, \texttt{worker2}) con \texttt{sudo NOPASSWD}, copia la clave pública de Vagrant a su \texttt{authorized\_keys} (para que el SSH desde el \emph{host} funcione con la misma clave de Vagrant) e instala una \emph{clave compartida del lab} en todas las VMs:

\begin{itemize}
  \item La clave privada Ed25519 se incrusta como \emph{heredoc} en el \emph{script}. Se autoriza también en \texttt{authorized\_keys} de cada usuario.
  \item Resultado: \texttt{master@test-master} puede hacer \texttt{ssh worker1@10.0.1.11} sin contraseña ni copia de claves manual.
\end{itemize}

Este mecanismo es el que después usan tanto Kubespray (Ansible necesita SSH \emph{passwordless} hacia los \emph{workers}) como \texttt{setup-lab.sh} (que entra por SSH a worker2 para hacer \texttt{chown} de los \texttt{hostPath}). El que la clave esté hardcodeada en el \emph{script} es aceptable porque el lab es local y la red no está expuesta.

\subsection{Trigger \texttt{ssh-config.sh}}

Se ejecuta en el \emph{host} tras levantar \texttt{test-worker2} (la última VM). Genera entradas SSH del estilo:

\begin{lstlisting}[style=bash]
Host test-master
  HostName 10.0.1.10
  User master
  Port 22
  IdentityFile /home/USER/.vagrant.d/insecure_private_key
\end{lstlisting}

dentro de \texttt{\textasciitilde/.ssh/config}, envueltas entre los marcadores \texttt{\# BEGIN paralelo-proyecto} / \texttt{\# END paralelo-proyecto}. La idempotencia se garantiza con \texttt{awk}: en cada ejecución borra el bloque anterior y lo reescribe. El usuario que aparece en \texttt{User} es el personalizado (\texttt{master}/\texttt{worker1}/\texttt{worker2}), no \texttt{vagrant}, mediante un \texttt{awk} que sustituye \texttt{User vagrant} por el correcto según el \texttt{Host}.

A partir de este punto el operador puede hacer \texttt{ssh test-master} desde el \emph{host} y entrar directamente como \texttt{master} con la clave de Vagrant.

\subsection{Discos persistentes}

\subsubsection{Decisión arquitectónica}

Los discos qcow2 \textbf{viven fuera de \texttt{vagrant-libvirt}}: son creados, formateados y montados por \emph{scripts} propios y no aparecen en el \texttt{Vagrantfile}. La razón es de durabilidad: \texttt{vagrant destroy -f} borra todo el almacenamiento que gestiona vagrant-libvirt (el disco principal de cada VM, los \emph{box} cacheados, etc.), pero los qcow2 de \texttt{disks/} sobreviven porque no son suyos. Esto permite reconstruir el clúster manteniendo los datos de \texttt{/mnt/data/grafana}, \texttt{/mnt/data/prometheus}, \texttt{/mnt/data/loki} entre destrucciones.

\subsubsection{\texttt{create-disks.sh} (\emph{trigger} en el host)}

Se ejecuta antes de levantar la primera VM, como \texttt{trigger.before :up only\_on: test-master}. Crea tres archivos qcow2 \emph{sparse} (la imagen pesa unos pocos MiB hasta que se escribe en el sistema de ficheros) dentro de \texttt{disks/}:

\begin{itemize}
  \item \texttt{test-master-data.qcow2}: 40\,GB
  \item \texttt{test-worker1-data.qcow2}: 20\,GB
  \item \texttt{test-worker2-data.qcow2}: 20\,GB
\end{itemize}

Antes de \texttt{qemu-img create}, el \emph{script} configura ACLs (\texttt{setfacl -m u:libvirt-qemu:x \$HOME}) para que el usuario \texttt{libvirt-qemu} pueda atravesar el \texttt{\$HOME} del usuario del \emph{host} y leer los qcow2. Esta ACL granular es preferible al \texttt{chmod o+x \$HOME}, que expondría todo el \texttt{home} al resto de usuarios del sistema.

\subsubsection{\texttt{attach-disks.sh} (\emph{trigger} en el host)}

Se ejecuta tras levantar \texttt{test-worker2}. Es la fase con más lógica idempotente del proyecto y trabaja en dos planos: \emph{host} (\texttt{virsh attach-disk}) y \emph{VM} (mediante SSH \emph{heredoc}). Estructura:

\begin{itemize}
  \item \textbf{Resolución del ID de dominio libvirt}: lee \texttt{.vagrant/machines/<node>/libvirt/id} para obtener el UUID del dominio. Sin ese fichero, omite el nodo con un mensaje de error.

  \item \textbf{Comprobación de \texttt{vdb} ya conectado}: \texttt{virsh domblklist} muestra los discos del dominio; si \texttt{vdb} ya aparece, se omite la conexión.

  \item \textbf{\texttt{virsh attach-disk \dots\ -{}-persistent}}: la opción \texttt{-{}-persistent} es clave. Sin ella, el \texttt{attach} se aplica solo al proceso QEMU vivo y se pierde en cualquier reinicio. Con \texttt{-{}-persistent}, libvirt escribe el cambio al XML de definición del dominio y \texttt{vdb} reaparece tras un \texttt{virsh start} o un reboot del \emph{host}.

  \item \textbf{Formateo dentro de la VM} (SSH \emph{heredoc}): si \texttt{blkid /dev/vdb} no encuentra sistema de ficheros, se formatea como \texttt{ext4}. Si ya existe, se respeta (idempotente).

  \item \textbf{Montaje persistente vía \texttt{/etc/fstab}}: se añade la línea \texttt{/dev/vdb /mnt/data ext4 defaults,nofail 0 2}. La opción \texttt{nofail} es defensiva: si por cualquier razón el disco no aparece, la VM arranca igualmente en vez de quedarse en modo \emph{emergency}.

  \item \textbf{Montaje en caliente}: \texttt{sudo mount /mnt/data} si no está ya montado.
\end{itemize}

El resultado es que cada nodo tiene un volumen ext4 estable bajo \texttt{/mnt/data}, listo para alojar \texttt{hostPath} de Kubernetes.

\subsection{Acceso gráfico a las VMs}

Las máquinas exponen consola gráfica VNC (driver \texttt{cirrus} por defecto en \texttt{generic/ubuntu2204}). Para acceder a una de ellas se puede usar:

\begin{lstlisting}[style=bash]
virt-viewer --connect qemu:///system test-master
# o desde virt-manager: conexión QEMU/KVM > test-master > Open Console
\end{lstlisting}

\subsection{Resumen del estado tras Fase 1}

Tras un \texttt{vagrant up} exitoso, el estado del laboratorio es:

\begin{itemize}
  \item Red \texttt{red-prueba} (\texttt{virbr2}) operativa con NAT y reserva DHCP por MAC.
  \item Tres VMs Ubuntu 22.04 con sus IPs estáticas (10.0.1.10/11/12).
  \item Usuarios personalizados \texttt{master}, \texttt{worker1}, \texttt{worker2} con \texttt{sudo NOPASSWD} y clave compartida del lab para SSH entre nodos.
  \item Acceso SSH desde el \emph{host} mediante alias (\texttt{ssh test-master}/\texttt{worker1}/\texttt{worker2}).
  \item Disco persistente qcow2 montado en \texttt{/mnt/data} en cada VM, listo para \texttt{hostPath} en fases posteriores.
\end{itemize}

A partir de aquí, los \emph{provisioners} \texttt{run: never} se invocan en el orden \texttt{setup-kubespray} $\to$ \texttt{setup-lab} $\to$ \texttt{setup-security}.
